Legal text is hard for named entity recognition. New statutes, parties and courts appear continually, and expert annotation is costly. Open-world NER assigns types that were not fixed at training time. These systems are built on general-domain text, though, and legal uses adapt them corpus by corpus and language by language. \rev{We take the opposite route, and apply OPENER unchanged}, a frugal pipeline of ready-to-use components. Its entity embedder is deliberately left generic, so its space is never fitted to one legal schema. We evaluate a zero-shot head (OPENER-ZS) and a supervised linear probe (OPENER-Sup) on four legal corpora, in English, Portuguese and German. The baselines are GLiNER, GNER, OWNER and a 4-bit Qwen, and we report quality (Adjusted Mutual Information), latency and energy jointly, over three seeds. The design holds under an English domain shift, and degrades cross-lingually. OPENER-ZS stays usable in English (up to 49.1 gold AMI), but it does not lead, and it weakens across languages (23.8 to 37.4). OPENER-Sup is barely affected (61.4 on average), and posts the best end-to-end score (35.9). Every system is then capped by the English detectors' recall, as low as 10.2\% on German. Both heads reach that ceiling at one to two orders of magnitude less latency and energy than the generative baselines. \rev{On detected mentions, the supervised head types 78.7\% correctly, against the detector's 64.6\%. A typing oracle over those same spans caps end-to-end quality at 70 AMI, the rest being mentions the detector never proposes.} Where labels exist, the supervised head is the practical operating point. Detection, not typing, is where adaptation should go first.
